\documentclass{article}
\usepackage{fontspec} 
\usepackage{unicode-math}  % 先加载 unicode-math
\setmathfont{XITS Math}    % 再设置数学字体
%\usepackage[UTF8]{ctex} % 如果需要支持中文

\usepackage{anyfontsize}
\usepackage{xeCJK} % XeLaTeX 推荐 xeCJK，LuaLaTeX 也可用   
\usepackage{setspace}  % 添加这个包
\setstretch{1.2}      % 设置行距为1.5倍

\setCJKmainfont[
    Path = C:/USERS/11252/APPDATA/LOCAL/MICROSOFT/WINDOWS/FONTS/,  % 改用 Windows 字体目录
    UprightFont = {SourceHanSansCN-Light1.otf},  % 使用可变字体
    BoldFont = {SourceHanSansCN-Medium1.otf},
    LetterSpace = 2.0,  % 增加字间距
    WordSpace = 1.2,    % 增加词间距
    %BoldFeatures = {RawFeature={+wght=900}},
    %Scale=1
]{SourceHanSansCN}


\setmainfont[
    Path = C:/USERS/11252/APPDATA/LOCAL/MICROSOFT/WINDOWS/FONTS/,  % 改用 Windows 字体目录
    UprightFont = {SourceHanSansCN-Light1.otf},  % 使用可变字体
    BoldFont = {SourceHanSansCN-Medium1.otf},
    LetterSpace = 2.0,  % 增加字间距
    WordSpace = 1.2,    % 增加词间距
    %BoldFeatures = {RawFeature={+wght=900}},
    %Scale=1
]{SourceHanSansCN}

%\setCJKmainfont[
 %   Path = C:/USERS/11252/APPDATA/LOCAL/MICROSOFT/WINDOWS/FONTS/,
 %   UprightFont = {SourceHanSerifCN-Light-5.otf},
 %   BoldFont = {SourceHanSerifCN-Medium-6.otf},
%    Scale=1
%]{SourceHanSerifCN}

%\setmainfont{SourceHanSerif}  % 设置英文字体

% 处理冲突的命令
\let\oldbackepsilon\backepsilon
\let\backepsilon\relax
\let\oldeth\eth
\let\eth\relax
\let\olddigamma\digamma
\let\digamma\relax
\usepackage{float}

\usepackage[a4paper, left=0.1in, right=0.1in, top=0.05in, bottom=0.05in]{geometry} % 设置四边页边距为0.1英寸{geometry} % 设置页边距为0.5英寸
\usepackage{enumitem} % 用于自定义列表
\usepackage{amsmath}  % 数学公式支持
\usepackage{tikz}
\usetikzlibrary{arrows.meta, positioning}
\setlength{\parindent}{0pt} % 不缩进
\usepackage{etoolbox} % 用于列表操作
%调整类代码
\usepackage{comment}
\usepackage{tocloft} % 用于定制目录 样式
\usepackage{hyperref} %不需要目录盒子注释这
\usepackage{ifthen}
\usepackage{tikz}
\usetikzlibrary{arrows.meta}
\usepackage{titletoc}
\usepackage{graphicx}
\usepackage{float}

\usepackage{amssymb}  % 提供数学符号，包括\therefore
%目录设定
%快捷键 CMD + / 注释
%  \renewcommand{\cftdot}{} % 隐藏目录中的页码
%  \cftpagenumbersoff{section}
%  \cftpagenumbersoff{subsection}
%  \makeatletter
%  \renewcommand{\@pnumwidth}{0em}  % 设置页码宽度为0
%  \renewcommand{\@tocrmarg}{0em}   % 设置右边距为0
%  \makeatother
 

\usepackage{environ}

\newif\ifshowcontent  % 定义一个条件判断变量
\showcontenttrue    % 设置为 false，隐藏所有 question 环境的内容

\newcounter{question} % 题目计数器
\setcounter{question}{0}
\NewEnviron{question}{%
    \ifshowcontent
        \par\stepcounter{question}\textbf{\thequestion.} \BODY\par % 如果显示内容，则输出
    \fi
}

\NewEnviron{choices}{%
    \ifshowcontent
        \begin{enumerate}[label=\Alph*., leftmargin=*]
            \BODY
        \end{enumerate}\par
    \fi
}

\NewEnviron{correctanswer}{%
    \ifshowcontent
        \textbf{答案:}\BODY\par
    \fi
}



\newenvironment{explanation}
{\textbf{解析:}\par}
{\par}



% 新增考察方面命令
\newcommand{\checklist}[1]{
    \textbf{考察:} #1 \par % 在题目下方显示考察方面
    \addcontentsline{toc}{subsection}{题号\thequestion 考察: #1} % 将考察内容添加到目录
    \vspace{2em} % 添加1em的垂直空间
}

\graphicspath{{images/}} %统一


\begin{document}
 %\fontsize{22}{31}\selectfont
 \tableofcontents % 添加目录
\newpage
%\pagestyle{empty}




\section{考点1:银行压力测试}
\setcounter{question}{0}

\begin{question}
    Which of the following statements regarding stress testing and firm wide risk aggregation is most accurate?
\end{question}

\begin{choices}
    \item One potential problem with the risk aggregation process is the underestimation of the diversification effect
    \item In the context of risk aggregation, stress tests are considered a more reliable alternative than regulatory and economic capital measures
    \item The process of stress testing can be performed in a relatively objective manner given that many common risks are reasonably quantifiable
    \item Using the results of stress testing, those scenarios that will result in minimal losses with a very low likelihood of occurrence may be ignored when considering any adjustments to risk appetite
\end{choices}

\begin{correctanswer}
    D
\end{correctanswer}

\begin{explanation}
    \textbf{A选项错误。}
    这个说法不准确：
    \begin{itemize}
        \item 风险加总常见问题是高估多样化效应
        \item 不是低估
    \end{itemize}

    \textbf{B选项错误。}
    这个说法不准确：
    \begin{itemize}
        \item 压力测试是对监管和经济资本计量的补充
        \item 不能完全替代
    \end{itemize}

    \textbf{C选项错误。}
    这个说法不准确：
    \begin{itemize}
        \item 压力测试需大量主观判断
        \item 很难做到完全客观
    \end{itemize}

    \textbf{D选项正确。}
    这个说法正确：
    \begin{itemize}
        \item 损失极小且概率极低的情景可在调整风险偏好时忽略
        \item 有助于聚焦主要风险情景
    \end{itemize}
\end{explanation}

\checklist{压力测试与风险加总的应用原则}


\begin{question}
    Nordlandia is a country with a developed economy maintaining its own currency, the Nordlandian crown (NLC), and whose most important export is domestically produced oil and natural gas. In a recent stress test of Nordlandia's banking system, several scenarios were considered. Which of the following is most consistent with being part of a coherent scenario?
\end{question}

\begin{choices}
    \item An increase in domestic inflation and appreciation of the NLC
    \item A significant increase in crude oil prices and a decrease in the Nordlandian housing price index
    \item A drop in crude oil prices and appreciation of the NLC
    \item A sustained decrease in natural gas prices and a decrease in the Nordlandian stock index
\end{choices}

\begin{correctanswer}
    D
\end{correctanswer}

\begin{explanation}
    \textbf{A选项错误。}
    这个说法不合理：
    \begin{itemize}
        \item 通胀上升通常导致货币贬值
        \item 与货币升值不一致
    \end{itemize}

    \textbf{B选项错误。}
    这个说法不合理：
    \begin{itemize}
        \item 油价大幅上涨通常带动经济和房价上涨
        \item 与房价下跌不一致
    \end{itemize}

    \textbf{C选项错误。}
    这个说法不合理：
    \begin{itemize}
        \item 油价下跌通常导致货币贬值
        \item 与货币升值不一致
    \end{itemize}

    \textbf{D选项正确。}
    这个说法合理：
    \begin{itemize}
        \item 天然气价格持续下跌削弱经济
        \item 股市下跌与经济疲软一致
        \item 场景逻辑连贯
    \end{itemize}
\end{explanation}

\checklist{压力测试场景设计的连贯性原则}


\begin{question}
    In a market crash the following are usually true? 
    I. Fixed-income portfolios hedged with short Treasury bonds and futures lose less than those hedged with interest rate swaps given equivalent durations. 
    II. Bid–offer spreads widen because of lower liquidity. 
    III. The spreads between off-the-run bonds and benchmark issues widen.
\end{question}

\begin{choices}
    \item I, II, and III
    \item II and III
    \item I and III
    \item None of the above
\end{choices}

\begin{correctanswer}
    B
\end{correctanswer}

\begin{explanation}
    \textbf{A选项错误。}
    这个说法不正确：
    \begin{itemize}
        \item I项错误：市场崩盘时，国债通常比利率互换涨幅更大
        \item 做空国债的损失可能大于做空互换
        \item II和III项正确
    \end{itemize}

    \textbf{B选项正确。}
    这个说法正确：
    \begin{itemize}
        \item II项：流动性下降导致买卖价差扩大
        \item III项：非主流债券与基准债券利差扩大
        \item 符合市场崩盘常见现象
    \end{itemize}

    \textbf{C选项错误。}
    这个说法不正确：
    \begin{itemize}
        \item I项错误，III项正确
        \item II项也应包括
    \end{itemize}

    \textbf{D选项错误。}
    这个说法不正确：
    \begin{itemize}
        \item II和III项在市场崩盘时普遍成立
    \end{itemize}
\end{explanation}

\checklist{市场崩盘时的流动性与利差变化}


\begin{question}
    Piper Hook, a bank examiner, is trying to make sense of stress tests done by one of the banks she examines. The stress tests are multi-factored and complex. The bank is using multiple extreme scenarios to test capital adequacy, making it difficult for Hook to interpret the results. One of the key stress test design challenges that Hook must deal with in her examination of stress tests is
\end{question}

\begin{choices}
    \item Multiplicity
    \item Efficiency
    \item Coherence
    \item Efficacy
\end{choices}

\begin{correctanswer}
    C
\end{correctanswer}

\begin{explanation}
    \textbf{A选项错误。}
    这个说法不准确：
    \begin{itemize}
        \item 多重性指情景数量多
        \item 但主要难点在于情景是否连贯
    \end{itemize}

    \textbf{B选项错误。}
    这个说法不准确：
    \begin{itemize}
        \item 效率不是压力测试设计的核心挑战
        \item 重点在于情景的合理性
    \end{itemize}

    \textbf{C选项正确。}
    这个说法正确：
    \begin{itemize}
        \item 连贯性是压力测试设计的关键难点
        \item 情景需极端但也要合理、可能
        \item 多因素情景更难保证连贯性
    \end{itemize}

    \textbf{D选项错误。}
    \begin{itemize}
        \item 有效性重要但不是主要设计难点
        \item 连贯性更为关键
    \end{itemize}
\end{explanation}

\checklist{压力测试设计中的连贯性原则}


\begin{question}
    Sarah Chen, a senior analyst at the Federal Reserve Bank, is presenting research on stress testing methodologies to a group of banking supervisors. She explains that macro-variable stress testing can be misleading for certain banks due to regional variations in economic indicators. She cites the example of significant differences in unemployment rates across European countries following the 2010-2012 sovereign debt crisis. Which type of loan did Chen most likely identify as having the strongest correlation with unemployment rates?
\end{question}

\begin{choices}
    \item Residential mortgage loans
    \item Credit card loans
    \item Commercial property loans
    \item Corporate term loans
\end{choices}

\begin{correctanswer}
    B
\end{correctanswer}

\begin{explanation}
    \textbf{A选项错误。}
    这个说法不准确：
    \begin{itemize}
        \item 住房抵押贷款虽受失业率影响
        \item 但违约敏感度低于信用卡贷款
    \end{itemize}

    \textbf{B选项正确。}
    这个说法正确：
    \begin{itemize}
        \item 信用卡损失对失业率变化最敏感
        \item 例如：希腊失业率达27\%，德国仅5\%
        \item 信用卡违约率地区差异显著
    \end{itemize}

    \textbf{C选项错误。}
    这个说法不准确：
    \begin{itemize}
        \item 商业地产贷款受经济周期影响更大
        \item 与失业率相关性相对较弱
    \end{itemize}

    \textbf{D选项错误。}
    这个说法不准确：
    \begin{itemize}
        \item 企业贷款违约主要受行业因素影响
        \item 与失业率相关性较低
    \end{itemize}
\end{explanation}

\checklist{压力测试中不同贷款类型对失业率的敏感度}





\newpage



\section{考点2:风险资本归属和风险调整后绩效衡量}
\setcounter{question}{0}
\begin{question}
    As part of its enterprise risk management enhancement initiative, Metropolitan Bank implemented a new economic capital-based decision framework that evaluates risk sources across various business divisions and organizational hierarchies. Which of the following statements about risk correlations is most accurate?
\end{question}

\begin{choices}
    \item Management decisions can influence the correlations between asset and liability risks
    \item Correlations between major risk categories (credit, market, operational) are well understood and easily measurable at the firm level
    \item Business unit correlations only matter for firm-wide economic capital allocation and are irrelevant for individual unit decisions
    \item Including correlations in firm-wide risk assessment typically results in a total VaR that is greater than or equal to the sum of individual unit VaRs
\end{choices}

\begin{correctanswer}
    A
\end{correctanswer}

\begin{explanation}
    \textbf{A选项正确。}
    这个说法正确：
    \begin{itemize}
        \item 管理层可通过调整资产负债结构影响风险相关性
        \item 例如：改变贷款组合或融资结构
        \item 这是风险管理的重要工具
    \end{itemize}

    \textbf{B选项错误。}
    这个说法不准确：
    \begin{itemize}
        \item 主要风险类别间的相关性难以准确估计
        \item 不同风险类型的数据和方法差异大
    \end{itemize}

    \textbf{C选项错误。}
    这个说法不准确：
    \begin{itemize}
        \item 业务单元相关性对单个单元决策也重要
        \item 影响资源分配和风险限额设定
    \end{itemize}

    \textbf{D选项错误。}
    这个说法不准确：
    \begin{itemize}
        \item 考虑相关性通常降低总VaR
        \item 因为存在风险分散效应
    \end{itemize}
\end{explanation}

\checklist{经济资本框架中风险相关性的管理}


\begin{question}
    Which of the following statements about economic and regulatory capital are correct?
    I. Regulatory capital aims to maintain banking system stability by ensuring adequate capital levels
    II. Economic capital is designed to maintain solvency at a specific confidence level
    III. For any bank, economic capital must always be lower than regulatory capital
    IV. Economic capital allocation is a strategic process that impacts business unit and overall bank performance
\end{question}

\begin{choices}
    \item II and IV only
    \item I, II, III, and IV
    \item I, II, and IV only
    \item I and IV only
\end{choices}

\begin{correctanswer}
    C
\end{correctanswer}

\begin{explanation}
    \textbf{A选项错误。}
    这个说法不完整：
    \begin{itemize}
        \item 遗漏了I项关于监管资本的正确描述
        \item II和IV项正确
    \end{itemize}

    \textbf{B选项错误。}
    这个说法不准确：
    \begin{itemize}
        \item III项错误：经济资本不一定低于监管资本
        \item 两者没有必然关系
    \end{itemize}

    \textbf{C选项正确。}
    这个说法正确：
    \begin{itemize}
        \item I项：监管资本确保银行体系稳定
        \item II项：经济资本维持特定置信水平下的偿付能力
        \item IV项：经济资本分配影响业务表现
        \item III项错误：经济资本与监管资本无必然大小关系
    \end{itemize}

    \textbf{D选项错误。}
    这个说法不完整：
    \begin{itemize}
        \item 遗漏了II项关于经济资本的正确描述
        \item I和IV项正确
    \end{itemize}
\end{explanation}

\checklist{经济资本与监管资本的区别与联系}


\begin{question}
    Your bank calculates a one-day 95\% VaR for market risk, a one-year 99\% VaR for operational risk, and a one-year 99\% VaR for credit risk. The measures are \$80 million, \$400 million, and \$800 million, respectively. Operational risk is defined to include all risks that are not market risks and credit risks, and these three categories are mutually uncorrelated. The market risk VaR assumes normally distributed returns, and the bank expects to maintain its market risk VaR at that level for the whole year. Your supervisor wants your best estimate of a firm-wide VaR at the 1\% level. Among the following choices, your best estimate is:
\end{question}

\begin{choices}
    \item \$1.36 billion
    \item \$1.55 billion
    \item \$2.00 billion
    \item It is impossible to aggregate risks with different distributions having only this information
\end{choices}

\begin{correctanswer}
    C
\end{correctanswer}

\begin{explanation}
    \textbf{A选项错误。}
    这个说法不准确：
    \begin{itemize}
        \item 未考虑市场风险VaR的时间调整
        \item 计算过程有误
    \end{itemize}

    \textbf{B选项错误。}
    这个说法不准确：
    \begin{itemize}
        \item 市场风险VaR的置信水平转换有误
        \item 未正确应用平方根法则
    \end{itemize}

    \textbf{C选项正确。}
    这个说法正确：
    \begin{itemize}
        \item 市场风险VaR调整：\$80 × (2.326/1.645) × √252 = \$1,796
        \item 总VaR = √(1,796² + 400² + 800²) = \$2,000
        \item 考虑了时间调整和置信水平转换
    \end{itemize}

    \textbf{D选项错误。}
    这个说法不准确：
    \begin{itemize}
        \item 有足够信息进行风险加总
        \item 已知分布假设和相关性
    \end{itemize}
\end{explanation}

\checklist{不同风险类型VaR的加总方法}


\begin{question}
    The Chairman of a commercial bank recognizes that inadequate cash flow planning is one of the most critical mistakes a financial institution can make and insists that a robust liquidity-at-risk (LaR) measurement system be a core component of the bank's risk management framework. Which of the following statements about LaR is most accurate?
\end{question}

\begin{choices}
    \item Hedging to reduce basis risk will lower LaR
    \item Futures hedging and long option positions have identical effects on LaR
    \item For a hedged portfolio, LaR can be substantially different from VaR
    \item A bank's LaR typically decreases as its credit rating deteriorates
\end{choices}

\begin{correctanswer}
    C
\end{correctanswer}

\begin{explanation}
    \textbf{A选项错误。}
    这个说法不准确：
    \begin{itemize}
        \item 对冲基差风险不一定降低LaR
        \item 可能增加保证金要求
    \end{itemize}

    \textbf{B选项错误。}
    这个说法不准确：
    \begin{itemize}
        \item 期货和期权对LaR影响不同
        \item 期货有保证金要求，期权有权利金
    \end{itemize}

    \textbf{C选项正确。}
    这个说法正确：
    \begin{itemize}
        \item 对冲组合的LaR与VaR可能显著不同
        \item 期货对冲可能增加保证金要求
        \item 欧式期权可能降低LaR
    \end{itemize}

    \textbf{D选项错误。}
    这个说法不准确：
    \begin{itemize}
        \item 信用质量下降通常增加LaR
        \item 融资成本上升，流动性压力增大
    \end{itemize}
\end{explanation}

\checklist{流动性风险与市场风险的关系}



\begin{question}
    You are a portfolio manager at a leading investment fund analyzing a 2,000 share position in an undervalued but illiquid stock XYZ, which has a current stock price of USD 85 (expressed as the midpoint of the current bid-ask spread.) Daily return for XYZ has an estimated volatility of 1.15%. The average bid-ask spread is USD 0.18. Assuming returns of XYZ are normally distributed, what is the estimated liquidity-adjusted daily 95% VaR, using the constant spread approach?
\end{question}

\begin{choices}
    \item USD 2,450
    \item USD 2,580
    \item USD 2,710
    \item USD 2,840
\end{choices}

\begin{correctanswer}
    C
\end{correctanswer}

\begin{explanation}
    \textbf{A选项错误。}
    这个说法不准确：
    \begin{itemize}
        \item 未正确计算流动性成本
        \item 计算过程有误
    \end{itemize}

    \textbf{B选项错误。}
    这个说法不准确：
    \begin{itemize}
        \item 仅计算了市场风险VaR
        \item 遗漏了流动性调整
    \end{itemize}

    \textbf{C选项正确。}
    这个说法正确：
    \begin{itemize}
        \item 市场风险VaR = 170,000 × (1.645 × 0.0115) = 2,580
        \item 流动性成本 = 170,000 × (0.5 × 0.18/85) = 130
        \item 流动性调整VaR = 2,580 + 130 = 2,710
    \end{itemize}

    \textbf{D选项错误。}
    这个说法不准确：
    \begin{itemize}
        \item 流动性成本计算过高
        \item 总VaR被高估
    \end{itemize}
\end{explanation}

\checklist{流动性调整VaR的计算方法}


\begin{question}
    Global Capital is an investment management firm with USD 30 billion under management. It holds 25% of the shares in a company. Global Capital's risk manager is concerned that, if the entire position needs to be liquidated, its size would impact the market price. The estimated price elasticity of demand is -0.4. What is the increase in Global Capital's Value-at-Risk estimate for this position if a liquidity adjustment is made?
\end{question}

\begin{choices}
    \item 5%
    \item 10%
    \item 15%
    \item 20%
\end{choices}

\begin{correctanswer}
    B
\end{correctanswer}

\begin{explanation}
    \textbf{A选项错误。}
    这个说法不准确：
    \begin{itemize}
        \item 未正确应用价格弹性公式
        \item 计算过程有误
    \end{itemize}

    \textbf{B选项正确。}
    这个说法正确：
    \begin{itemize}
        \item 价格变化率 = 弹性 × 交易规模比例
        \item ΔP/P = -0.4 × 0.25 = -0.1
        \item LVaR/VaR = 1 - (-0.1) = 1.1
        \item VaR增加10%
    \end{itemize}

    \textbf{C选项错误。}
    这个说法不准确：
    \begin{itemize}
        \item 高估了流动性影响
        \item 计算过程有误
    \end{itemize}

    \textbf{D选项错误。}
    这个说法不准确：
    \begin{itemize}
        \item 显著高估了流动性调整
        \item 不符合价格弹性计算
    \end{itemize}
\end{explanation}

\checklist{流动性调整对VaR的影响计算}


\begin{question}
    In calculating its risk-adjusted return on capital, your bank uses a capital charge of 3.00% for revolving credit facilities with a loan equivalent factor of 0.40 assigned to the undrawn portion. Recently, you have become concerned that the protective covenants in these loans are insufficient and may not prevent customers from drawing on the facilities during stress periods. As such, you have recommended increasing the loan equivalent factor to 0.80. This recommendation has met with resistance from the loan origination team, and senior management has asked you to quantify the impact of your recommendation. For a typical facility that has an original principal of USD 1.2 billion and is 25% drawn, how much additional economic capital would have to be allocated if you increase the loan equivalent factor from 0.40 to 0.80?
\end{question}

\begin{choices}
    \item USD 4.32 million
    \item USD 7.20 million
    \item USD 10.80 million
    \item USD 15.12 million
\end{choices}

\begin{correctanswer}
    B
\end{correctanswer}

\begin{explanation}
    \textbf{A选项错误。}
    这个说法不准确：
    \begin{itemize}
        \item 未正确计算未提用部分的资本要求
        \item 计算过程有误
    \end{itemize}

    \textbf{B选项正确。}
    这个说法正确：
    \begin{itemize}
        \item 初始资本要求 = [(12亿 × 0.25) + (12亿 × 0.75 × 0.40)] × 3% = 16.2百万
        \item 调整后资本要求 = [(12亿 × 0.25) + (12亿 × 0.75 × 0.80)] × 3% = 23.4百万
        \item 额外资本要求 = 23.4 - 16.2 = 7.2百万
    \end{itemize}

    \textbf{C选项错误。}
    这个说法不准确：
    \begin{itemize}
        \item 高估了资本要求变化
        \item 计算过程有误
    \end{itemize}

    \textbf{D选项错误。}
    这个说法不准确：
    \begin{itemize}
        \item 显著高估了资本要求
        \item 未正确应用公式
    \end{itemize}
\end{explanation}

\checklist{循环信贷额度经济资本计算}


\begin{question}
    Michael Chen, an analyst at Global Advisors, is analyzing the economic effects of buying stock with borrowed funds for a high net worth individual client. Assume that the client has \$300 cash invested (i.e., no borrowed funds) and then uses the cash to purchase stock. The client then decides to use 60\% borrowed funds to purchase stock on margin. After the margin transaction, the total assets on the full economic balance sheet and the leverage ratio are closest to:

    A. Total assets: \$300; Leverage ratio: 1.6
    B. Total assets: \$450; Leverage ratio: 1.6
    C. Total assets: \$450; Leverage ratio: 2.5
    D. Total assets: \$600; Leverage ratio: 2.5
\end{question}

\begin{choices}
    \item A
    \item B
    \item C
    \item D
\end{choices}

\begin{correctanswer}
    B
\end{correctanswer}

\begin{explanation}
    \textbf{A选项错误。}
    这个说法不准确：
    \begin{itemize}
        \item 未考虑保证金交易后的总资产变化
        \item 杠杆率计算有误
    \end{itemize}

    \textbf{B选项正确。}
    这个说法正确：
    \begin{itemize}
        \item 初始状态：现金\$300，无负债
        \item 保证金交易：借入\$180（60%）
        \item 总资产 = \$300 + \$180 = \$450
        \item 杠杆率 = 总资产/权益 = \$450/\$300 = 1.6
    \end{itemize}

    \textbf{C选项错误。}
    这个说法不准确：
    \begin{itemize}
        \item 杠杆率计算错误
        \item 应为1.6而非2.5
    \end{itemize}

    \textbf{D选项错误。}
    这个说法不准确：
    \begin{itemize}
        \item 总资产和杠杆率都计算错误
        \item 高估了保证金交易的影响
    \end{itemize}
\end{explanation}

\checklist{保证金交易的杠杆效应计算}


\begin{question}
    A trader observes a quote for Stock ABC, and the midpoint of its current best bid and best ask price is CAD 42. ABC has an estimated daily return volatility of 0.30% and average bid-ask spread of CAD 0.12. Assuming the returns of ABC are normally distributed, what is closest to the estimated liquidity-adjusted, 1-day 95% VaR, using the constant spread approach on a 15,000 share position?
\end{question}

\begin{choices}
    \item CAD 2,500
    \item CAD 3,200
    \item CAD 3,900
    \item CAD 4,600
\end{choices}

\begin{correctanswer}
    B
\end{correctanswer}

\begin{explanation}
    \textbf{A选项错误。}
    这个说法不准确：
    \begin{itemize}
        \item 仅计算了市场风险VaR
        \item 遗漏了流动性成本
    \end{itemize}

    \textbf{B选项正确。}
    这个说法正确：
    \begin{itemize}
        \item 市场风险VaR = 42 × 15,000 × (1.645 × 0.003) = 2,700
        \item 流动性成本 = 630,000 × (0.5 × 0.12/42) = 500
        \item 流动性调整VaR = 2,700 + 500 = 3,200
    \end{itemize}

    \textbf{C选项错误。}
    这个说法不准确：
    \begin{itemize}
        \item 高估了流动性成本
        \item 计算过程有误
    \end{itemize}

    \textbf{D选项错误。}
    这个说法不准确：
    \begin{itemize}
        \item 显著高估了总VaR
        \item 未正确应用公式
    \end{itemize}
\end{explanation}

\checklist{流动性调整VaR的计算方法}


\begin{question}
    A firm has determined that the risk-adjusted return on capital (RAROC) for a particular project is 16%. To evaluate whether the firm should accept the project, an analyst determines that the firm's beta is 1.4, the expected market return is 14%, and the risk-free interest rate is 6%. If the analyst uses the adjusted RAROC (ARAROC) methodology to make an accept/reject decision, should the project be accepted?
\end{question}

\begin{choices}
    \item No, because the computed ARAROC is approximately 1% less than the market risk premium
    \item No, because the RAROC is 1.25% less than the return predicted by the CAPM
    \item Yes, because the computed ARAROC is approximately 1% more than the market risk premium
    \item Yes, because the ARAROC is approximately 4% more than the return predicted by the CAPM
\end{choices}

\begin{correctanswer}
    A
\end{correctanswer}

\begin{explanation}
    \textbf{A选项正确。}
    这个说法正确：
    \begin{itemize}
        \item ARAROC = (RAROC - Rf)/β = (0.16 - 0.06)/1.4 = 7.14%
        \item 市场风险溢价 = 14% - 6% = 8%
        \item ARAROC比市场风险溢价低约1%
    \end{itemize}

    \textbf{B选项错误。}
    这个说法不准确：
    \begin{itemize}
        \item 未正确计算ARAROC
        \item 比较基准有误
    \end{itemize}

    \textbf{C选项错误。}
    这个说法不准确：
    \begin{itemize}
        \item ARAROC低于市场风险溢价
        \item 不应接受项目
    \end{itemize}

    \textbf{D选项错误。}
    这个说法不准确：
    \begin{itemize}
        \item 高估了ARAROC与CAPM预测收益的差异
        \item 计算过程有误
    \end{itemize}
\end{explanation}

\checklist{ARAROC与市场风险溢价的比较}


\begin{question}
    A bank issues a \$250,000,000 loan with the following characteristics:
    • Loan pays a fixed annual interest rate of 9.0%
    • The interest expense associated with the loan is 6.5%
    • The operating cost to the bank's commercial lending division is \$2,000,000
    • Economic capital required to support the loan is \$20 million, which is invested in T-bills paying a rate of 3.0\%
    • The expected loss for the loan is 20 basis points per year
    What is the risk-adjusted return on capital (RAROC) for this loan?
\end{question}

\begin{choices}
    \item 3.50\%
    \item 7.25\%
    \item 19.50\%
    \item 23.50\%
\end{choices}

\begin{correctanswer}
    D
\end{correctanswer}

\begin{explanation}
    \textbf{A选项错误。}
    这个说法不准确：
    \begin{itemize}
        \item 未考虑所有收入和成本项目
        \item 计算过程有误
    \end{itemize}

    \textbf{B选项错误。}
    这个说法不准确：
    \begin{itemize}
        \item 遗漏了经济资本收益
        \item 未正确计算RAROC
    \end{itemize}

    \textbf{C选项错误。}
    这个说法不准确：
    \begin{itemize}
        \item 高估了预期损失影响
        \item 计算过程有误
    \end{itemize}

    \textbf{D选项正确。}
    这个说法正确：
    \begin{itemize}
        \item 收入 = 250,000,000 × 9\% = 22,500,000
        \item 预期损失 = 250,000,000 × 0.2\% = 500,000
        \item 费用 = 250,000,000 × 6.5\% + 2,000,000 = 18,250,000
        \item 经济资本收益 = 20,000,000 × 3\% = 600,000
        \item RAROC = (22,500,000 - 500,000 - 18,250,000 + 600,000)/20,000,000 = 23.5\%
    \end{itemize}
\end{explanation}

\checklist{RAROC的计算方法}


\begin{question}
    Suppose a bank loan has the following characteristics:

    • Gross revenues are \$1.5 million

    • Interest expense is \$0.4 million

    • The return on the \$6 million economic capital invested in T-bills is \$150,000

    • The firm beta is 1.6

    • The unexpected loss for the loan is estimated at \$600,000

    • The adjusted RAROC is equal to 9\%

    What is the worst-case loss for this loan?
\end{question}

\begin{choices}
    \item \$1.8 million
    \item \$2.5 million
    \item \$1.1 million
    \item \$0.85 million
\end{choices}

\begin{correctanswer}
    C
\end{correctanswer}

\begin{explanation}
    \textbf{A选项错误。}
    这个说法不准确：
    \begin{itemize}
        \item 高估了最坏情况损失
        \item 计算过程有误
    \end{itemize}

    \textbf{B选项错误。}
    这个说法不准确：
    \begin{itemize}
        \item 显著高估了损失
        \item 未正确应用公式
    \end{itemize}

    \textbf{C选项正确。}
    这个说法正确：
    \begin{itemize}
        \item 无风险利率 = 150,000/6,000,000 = 2.5\%
        \item RAROC = 1.6 × 0.09 + 0.025 = 16.9\%
        \item 预期损失 = -(0.169 × 6M - 1.5M + 0.4M - 0.15M) = 500,000
        \item 最坏情况损失 = 500,000 + 600,000 = 1.1M
    \end{itemize}

    \textbf{D选项错误。}
    这个说法不准确：
    \begin{itemize}
        \item 低估了最坏情况损失
        \item 未正确计算预期损失
    \end{itemize}
\end{explanation}

\checklist{最坏情况损失的计算方法}

\begin{question}
    In its efforts to enhance its enterprise risk management function, Metropolitan Bank implemented a new economic capital-based decision framework that evaluates risk sources across various business divisions and organizational hierarchies. Which of the following statements about risk correlations is most accurate?
\end{question}

\begin{choices}
    \item Management decisions can influence the correlations between asset and liability risks
    \item Correlations between major risk categories (credit, market, operational) are well understood and easily measurable at the firm level
    \item Business unit correlations only matter for firm-wide economic capital allocation and are irrelevant for individual unit decisions
    \item Including correlations in firm-wide risk assessment typically results in a total VaR that is greater than or equal to the sum of individual unit VaRs
\end{choices}

\begin{correctanswer}
    A
\end{correctanswer}

\begin{explanation}
    \textbf{A选项正确。}
    这个说法正确：
    \begin{itemize}
        \item 管理层可通过调整资产负债结构影响风险相关性
        \item 例如：改变贷款组合或融资结构
        \item 这是风险管理的重要工具
    \end{itemize}

    \textbf{B选项错误。}
    这个说法不准确：
    \begin{itemize}
        \item 主要风险类别间的相关性难以准确估计
        \item 不同风险类型的数据和方法差异大
    \end{itemize}

    \textbf{C选项错误。}
    这个说法不准确：
    \begin{itemize}
        \item 业务单元相关性对单个单元决策也重要
        \item 影响资源分配和风险限额设定
    \end{itemize}

    \textbf{D选项错误。}
    这个说法不准确：
    \begin{itemize}
        \item 考虑相关性通常降低总VaR
        \item 因为存在风险分散效应
    \end{itemize}
\end{explanation}

\checklist{经济资本框架中风险相关性的管理}


\begin{question}
    A treasury analyst at a bank is using the RAROC approach to assess whether to extend a new CNY 1 billion loan. The loan would be fully funded by deposits and would have the following metrics:
    • Required economic capital: CNY 120 million
    • Interest rate on the loan: 5.5\%
    • Operating expenses to support the loan: CNY 8 million
    • Expected loss on the loan: CNY 15 million
    • Unexpected loss on the loan: CNY 12 million
    • Transfer price charge on deposits funding the loan: 1.2\%
    • Effective tax rate: 25\%
    The analyst assumes a return of 2.5\% on the bank's economic capital over the duration of the loan. What is the after-tax RAROC for this loan?
\end{question}

\begin{choices}
    \item 5.5\%
    \item 6.2\%
    \item 10.8\%
    \item 12.5\%
\end{choices}

\begin{correctanswer}
    D
\end{correctanswer}

\begin{explanation}
    \textbf{A选项错误。}
    这个说法不准确：
    \begin{itemize}
        \item 未考虑所有收入和成本项目
        \item 计算过程有误
    \end{itemize}

    \textbf{B选项错误。}
    这个说法不准确：
    \begin{itemize}
        \item 遗漏了经济资本收益
        \item 未正确计算RAROC
    \end{itemize}

    \textbf{C选项错误。}
    这个说法不准确：
    \begin{itemize}
        \item 高估了预期损失影响
        \item 计算过程有误
    \end{itemize}

    \textbf{D选项正确。}
    这个说法正确：
    \begin{itemize}
        \item 收入 = 10亿 × 5.5\% = 5,500万
        \item 费用 = 10亿 × 1.2\% + 800万 = 2,000万
        \item 预期损失 = 1,500万
        \item 经济资本收益 = 1.2亿 × 2.5\% = 300万
        \item 税前收益 = 5,500万 - 2,000万 - 1,500万 + 300万 = 2,300万
        \item 税后RAROC = 2,300万 × (1-25\%)/1.2亿 = 12.5\%
    \end{itemize}
\end{explanation}

\checklist{税后RAROC的计算方法}





\newpage


\section{考点3:经济资本框架中的各种做法和问题}
\setcounter{question}{0}
\begin{question}
    Which of the following statements about the differences between market and operational value-at-risk at financial institutions are correct?
    I. The distribution of operational risk events must include sufficient mass in the extreme tail, making an assumption of a normal distribution invalid
    II. The typical time horizon of market VaR calculations is 1 day, whereas the typical time horizon of operational VaR calculations is 1 year
    III. Since prices are sufficiently available for liquid assets at all times, the market risk of liquid assets can be modeled using continuous distributions, but the nature of operational risk events requires using discrete distributions
    IV. Market VaR requires a higher confidence level than operational VaR
\end{question}

\begin{choices}
    \item I, II, and III
    \item I, II and IV
    \item I, II, III and IV
    \item III and IV
\end{choices}

\begin{correctanswer}
    A
\end{correctanswer}

\begin{explanation}
    \textbf{A选项正确。}
    这个说法正确：
    \begin{itemize}
        \item I项：操作风险分布尾部特征明显，不能用正态分布
        \item II项：市场风险VaR通常为1天，操作风险VaR通常为1年
        \item III项：市场风险可用连续分布，操作风险需用离散分布
    \end{itemize}

    \textbf{B选项错误。}
    这个说法不准确：
    \begin{itemize}
        \item IV项错误：市场风险VaR置信水平通常低于操作风险
        \item 置信水平是用户设定的参数
    \end{itemize}

    \textbf{C选项错误。}
    这个说法不准确：
    \begin{itemize}
        \item 包含IV项错误说法
        \item 其他三项正确
    \end{itemize}

    \textbf{D选项错误。}
    这个说法不准确：
    \begin{itemize}
        \item 仅III项正确
        \item IV项错误
        \item 遗漏I和II项正确说法
    \end{itemize}
\end{explanation}

\checklist{市场风险与操作风险VaR的区别}



\begin{question}
    Your Chief Risk Officer asks you to prepare a list of early warning indicators for liquidity problems for your bank. Which of the following are early warning indicators of a potential liquidity problem?
    I. Rapid asset growth, especially when funded with potentially volatile liabilities
    II. Growing concentrations in assets or liabilities
    III. An increase of the weighted average maturity of liabilities
    IV. Reduction in the frequency of positions approaching or breaching internal or regulatory limits
    V. Narrowing debt or credit default swap spreads
    VI. Counterparties that request additional collateral for credit exposures
    VII. Increasing redemptions of CDs before maturity
\end{question}

\begin{choices}
    \item I, II, VI, and VII
    \item I, III, V, and VI
    \item II, IV, V, and VII
    \item I, V, VI, and VII
\end{choices}

\begin{correctanswer}
    A
\end{correctanswer}

\begin{explanation}
    \textbf{A选项正确。}
    这个说法正确：
    \begin{itemize}
        \item I项：快速资产增长，特别是依赖不稳定负债融资
        \item II项：资产或负债集中度增加
        \item VI项：交易对手要求增加抵押品
        \item VII项：定期存款提前赎回增加
    \end{itemize}

    \textbf{B选项错误。}
    这个说法不准确：
    \begin{itemize}
        \item III项：负债期限延长降低短期融资风险
        \item V项：利差收窄不是流动性问题指标
    \end{itemize}

    \textbf{C选项错误。}
    这个说法不准确：
    \begin{itemize}
        \item IV项：这是市场风险指标
        \item V项：利差收窄不是问题指标
    \end{itemize}

    \textbf{D选项错误。}
    这个说法不准确：
    \begin{itemize}
        \item V项：利差收窄不是问题指标
        \item 其他三项正确
    \end{itemize}
\end{explanation}

\checklist{流动性风险的早期预警指标}


\begin{question}
    Metropolitan Bank approaches economic capital and risk aggregation by first estimating the stand-alone economic capital for individual risk factors. In a second step, the bank aggregates risks based on the relative amounts of economic capital allocated to these risks, taking into account the correlations between risk factors. Which of the following variables is not a primary driver of the diversification benefit that accrues from aggregation?
\end{question}

\begin{choices}
    \item The number of risk positions
    \item The size of the portfolio
    \item The concentration of those risk positions, or their relative weights in a portfolio
    \item The correlation between the positions
\end{choices}

\begin{correctanswer}
    B
\end{correctanswer}

\begin{explanation}
    \textbf{A选项错误。}
    这个说法不准确：
    \begin{itemize}
        \item 风险头寸数量是分散化收益的主要驱动因素
        \item 头寸数量越多，分散化效果越好
    \end{itemize}

    \textbf{B选项正确。}
    这个说法正确：
    \begin{itemize}
        \item 组合规模不是分散化收益的主要驱动因素
        \item 风险度量与组合规模成正比
        \item 规模翻倍，风险也翻倍
    \end{itemize}

    \textbf{C选项错误。}
    这个说法不准确：
    \begin{itemize}
        \item 风险头寸集中度是分散化收益的主要驱动因素
        \item 集中度越低，分散化效果越好
    \end{itemize}

    \textbf{D选项错误。}
    这个说法不准确：
    \begin{itemize}
        \item 头寸间相关性是分散化收益的主要驱动因素
        \item 相关性越低，分散化效果越好
    \end{itemize}
\end{explanation}

\checklist{风险分散化收益的驱动因素}


\begin{question}
    Consider a bank that wants to have an amount of capital so that it can absorb unexpected losses corresponding to a firm-wide VaR at the 1% level. It measures firm-wide VaR by adding up the VaRs for market risk, operational risk, and credit risk. There is a risk that the bank has too little capital because
\end{question}

\begin{choices}
    \item It does not take into account the correlations among risks
    \item It ignores risks that are not market, operational, or credit risks
    \item It mistakenly uses VaR to measure operational risk because operational risks that matter are rare events
    \item It is meaningless to add VaRs
\end{choices}

\begin{correctanswer}
    B
\end{correctanswer}

\begin{explanation}
    \textbf{A选项错误。}
    这个说法不准确：
    \begin{itemize}
        \item 忽略相关性会导致资本要求过高
        \item 而不是资本不足
    \end{itemize}

    \textbf{B选项正确。}
    这个说法正确：
    \begin{itemize}
        \item 仅考虑三大风险类别可能低估总风险
        \item 其他风险类型（如流动性风险）未被纳入
        \item 导致资本要求不足
    \end{itemize}

    \textbf{C选项错误。}
    这个说法不准确：
    \begin{itemize}
        \item 操作风险VaR可以包含罕见事件
        \item 不是资本不足的原因
    \end{itemize}

    \textbf{D选项错误。}
    这个说法不准确：
    \begin{itemize}
        \item VaR可以跨风险类型加总
        \item 加总会高估资本要求
    \end{itemize}
\end{explanation}

\checklist{银行资本充足性的风险覆盖范围}


\begin{question}
    Based on "Supervisory Guidance for Assessing Banks' Financial Instrument Fair Value Practices" issued by the Basel Committee, which of the following factors should be considered in determining whether the sources of fair values are reliable and relevant?

    i. Frequency and availability of prices/quotes
    ii. Maturity of the market

    iii. Agreement of values with those generated by internal models

    iv. Number of independent sources that produce the prices/quotes
\end{question}

\begin{choices}
    \item i and ii only
    \item iii and iv only
    \item i, ii and iii only
    \item i, ii and iv only
\end{choices}

\begin{correctanswer}
    D
\end{correctanswer}

\begin{explanation}
    \textbf{A选项错误。}
    这个说法不完整：
    \begin{itemize}
        \item 遗漏了独立来源数量
        \item i和ii项正确
    \end{itemize}

    \textbf{B选项错误。}
    这个说法不准确：
    \begin{itemize}
        \item iii项错误：内部模型一致性不是必要因素
        \item iv项正确
    \end{itemize}

    \textbf{C选项错误。}
    这个说法不准确：
    \begin{itemize}
        \item iii项错误：内部模型一致性不是必要因素
        \item i和ii项正确
    \end{itemize}

    \textbf{D选项正确。}
    这个说法正确：
    \begin{itemize}
        \item i项：价格/报价的频率和可得性
        \item ii项：市场成熟度
        \item iv项：独立来源数量
        \item iii项不是必要因素
    \end{itemize}
\end{explanation}

\checklist{公允价值可靠性的评估因素}


\begin{question}
    Alpha Capital, LLC (Alpha) uses monthly model vetting to mitigate potential model risk. Alpha's managers recently accepted the use of a model for valuing short-term options on 25-year corporate bonds, but rejected the same model to value short-term options on five-year government bonds. The managers also frequently test proposed analytical models against a simulation approach. These model vetting techniques are examples of which of the following vetting phases?
    Accepting/rejecting a model
    Testing models against simulation
\end{question}

\begin{choices}
    \item Health check of the model Stress testing
    \item Soundness of a model Stress testing
    \item Health check of the model Benchmark modeling
    \item Soundness of a model Benchmark modeling
\end{choices}

\begin{correctanswer}
    D
\end{correctanswer}

\begin{explanation}
    \textbf{A选项错误。}
    这个说法不准确：
    \begin{itemize}
        \item 健康检查确保模型包含必要属性
        \item 压力测试检查模型对不同情况的反应
    \end{itemize}

    \textbf{B选项错误。}
    这个说法不准确：
    \begin{itemize}
        \item 模型健全性检验正确
        \item 但压力测试不是模拟测试的正确描述
    \end{itemize}

    \textbf{C选项错误。}
    这个说法不准确：
    \begin{itemize}
        \item 健康检查不是模型接受/拒绝的正确描述
        \item 基准建模正确
    \end{itemize}

    \textbf{D选项正确。}
    这个说法正确：
    \begin{itemize}
        \item 模型健全性：确保模型适用于特定资产
        \item 基准建模：用模拟方法测试分析模型
    \end{itemize}
\end{explanation}

\checklist{模型验证的不同阶段}


\begin{question}
    Which of the following risk aggregation methodologies is characterized by great difficulty in validating parameterization and building a joint distribution?
\end{question}

\begin{choices}
    \item Copulas
    \item Constant diversification
    \item Variance-covariance matrix
    \item Full modeling/Simulation
\end{choices}

\begin{correctanswer}
    A
\end{correctanswer}

\begin{explanation}
    \textbf{A选项正确。}
    这个说法正确：
    \begin{itemize}
        \item Copula方法有两个主要缺点
        \item 参数验证非常困难
        \item 构建联合分布非常复杂
    \end{itemize}

    \textbf{B选项错误。}
    这个说法不准确：
    \begin{itemize}
        \item 固定分散化方法相对简单
        \item 不需要复杂的参数验证
    \end{itemize}

    \textbf{C选项错误。}
    这个说法不准确：
    \begin{itemize}
        \item 方差-协方差矩阵方法较为直观
        \item 参数验证相对容易
    \end{itemize}

    \textbf{D选项错误。}
    这个说法不准确：
    \begin{itemize}
        \item 全模型/模拟方法参数验证较直接
        \item 联合分布构建相对简单
    \end{itemize}
\end{explanation}

\checklist{风险加总方法的参数验证难度}


\begin{question}
    Which of the following model validation processes is specifically characterized by the limitation that it provides little comfort that the model actually reflects reality?
\end{question}

\begin{choices}
    \item Backtesting
    \item Benchmarking
    \item Stress testing
    \item Qualitative review
\end{choices}

\begin{correctanswer}
    B
\end{correctanswer}

\begin{explanation}
    \textbf{A选项错误。}
    这个说法不准确：
    \begin{itemize}
        \item 回测直接比较模型预测与实际结果
        \item 能较好反映模型对现实的拟合程度
    \end{itemize}

    \textbf{B选项正确。}
    这个说法正确：
    \begin{itemize}
        \item 基准测试只能比较不同模型
        \item 无法确认模型是否反映现实
        \item 仅能验证参数或输出的可比性
    \end{itemize}

    \textbf{C选项错误。}
    这个说法不准确：
    \begin{itemize}
        \item 压力测试检验模型在极端情况下的表现
        \item 能反映模型对现实风险的捕捉
    \end{itemize}

    \textbf{D选项错误。}
    这个说法不准确：
    \begin{itemize}
        \item 定性审查评估模型假设和逻辑
        \item 能验证模型对现实的反映
    \end{itemize}
\end{explanation}

\checklist{模型验证方法的局限性}



\begin{question}
    A bank credit officer, who has reviewed a loan application, has made the following statement: "On a standalone basis, I was not very keen on granting this loan; however, I granted this loan after looking at the overall asset portfolio of the bank." Based on the above statement, which of the following is true?
\end{question}

\begin{choices}
    \item The correlation of the newly granted loan with the overall portfolio is low and therefore the credit officer was right in granting the loan
    \item The correlation of the newly granted loan with the overall portfolio is low and therefore the credit officer was wrong in granting the loan
    \item The correlation of the newly granted loan with the overall portfolio is high and therefore the credit officer was right in granting the loan
    \item The correlation of the newly granted loan with the overall portfolio is high and therefore the credit officer was wrong in granting the loan
\end{choices}

\begin{correctanswer}
    A
\end{correctanswer}

\begin{explanation}
    \textbf{A选项正确。}
    这个说法正确：
    \begin{itemize}
        \item 贷款风险由系统性和非系统性风险共同决定
        \item 低相关性带来分散化收益
        \item 信贷员的决策合理
    \end{itemize}

    \textbf{B选项错误。}
    这个说法不准确：
    \begin{itemize}
        \item 低相关性有利于分散风险
        \item 不应拒绝贷款
    \end{itemize}

    \textbf{C选项错误。}
    这个说法不准确：
    \begin{itemize}
        \item 高相关性增加组合风险
        \item 不应批准贷款
    \end{itemize}

    \textbf{D选项错误。}
    这个说法不准确：
    \begin{itemize}
        \item 高相关性确实不利
        \item 但信贷员已批准贷款
    \end{itemize}
\end{explanation}

\checklist{贷款组合风险分散化原则}


\begin{question}
    Which of the following is not a major challenge for banks in reporting measures of conduct and culture?
\end{question}

\begin{choices}
    \item Quality data are often not available to enable boards and management to drill down into detail at the business level
    \item Standardized measures are more difficult to define for smaller regional banks than for global banks
    \item Limited, asymmetric data make comparisons difficult for identifying trends and relationships with other variables
    \item Management teams at banks are inexperienced in implementing a reporting process for conduct and culture
\end{choices}

\begin{correctanswer}
    B
\end{correctanswer}

\begin{explanation}
    \textbf{A选项错误。}
    这个说法不准确：
    \begin{itemize}
        \item 数据质量不足是主要挑战
        \item 影响管理层深入分析
    \end{itemize}

    \textbf{B选项正确。}
    这个说法正确：
    \begin{itemize}
        \item 标准化指标定义困难是普遍问题
        \item 与银行规模无关
        \item 主要挑战在于跨业务和地区的一致性
    \end{itemize}

    \textbf{C选项错误。}
    这个说法不准确：
    \begin{itemize}
        \item 数据有限且不对称是主要挑战
        \item 影响趋势分析和变量关系识别
    \end{itemize}

    \textbf{D选项错误。}
    这个说法不准确：
    \begin{itemize}
        \item 管理层经验不足是主要挑战
        \item 影响报告流程实施
    \end{itemize}
\end{explanation}

\checklist{银行行为与文化报告的主要挑战}

\begin{question}
    The use of which of the following items is meant more for protecting against risk deterioration?
\end{question}

\begin{choices}
    \item Risk based pricing
    \item Management incentives
    \item Credit portfolio management
    \item Customer profitability analysis
\end{choices}

\begin{correctanswer}
    C
\end{correctanswer}

\begin{explanation}
    \textbf{A选项错误。}
    这个说法不准确：
    \begin{itemize}
        \item 风险定价主要用于最大化银行利润
        \item 不是防范风险恶化的工具
    \end{itemize}

    \textbf{B选项错误。}
    这个说法不准确：
    \begin{itemize}
        \item 管理激励用于鼓励管理者参与经济资本配置
        \item 不是防范风险恶化的工具
    \end{itemize}

    \textbf{C选项正确。}
    这个说法正确：
    \begin{itemize}
        \item 信用组合管理是防范风险恶化的主要工具
        \item 通过组合管理控制整体风险水平
    \end{itemize}

    \textbf{D选项错误。}
    这个说法不准确：
    \begin{itemize}
        \item 客户盈利分析用于识别不盈利客户
        \item 不是防范风险恶化的工具
    \end{itemize}
\end{explanation}

\checklist{风险恶化防范工具的选择}



\begin{question}
    A risk manager at bank ABC is evaluating approaches the bank can use to measure the total risk in its portfolio. The bank's portfolio contains foreign currency forwards and swaps, credit instruments and derivatives, as well as various types of loans. The manager is studying how the compartmentalized and the unified risk measurement approaches are applied and the differences between them. Which of the following is correct?
\end{question}

\begin{choices}
    \item If diversification benefits are present in the portfolio, the compartmentalized approach will typically generate a lower estimate of the bank's total risk than the unified approach
    \item Interactions between risk categories that result in compounding effects are more appropriately captured by using the compartmentalized approach rather than the unified approach
    \item The unified approach sums the amount of risk separately measured in each risk category into a single estimate of the bank's total risk
    \item All else held constant, if wrong-way risk is present in the portfolio, the unified approach will typically generate a higher estimate of the bank's total risk than the compartmentalized approach
\end{choices}

\begin{correctanswer}
    D
\end{correctanswer}

\begin{explanation}
    \textbf{A选项错误。}
    这个说法不准确：
    \begin{itemize}
        \item 分散化收益在统一方法中更明显
        \item 分块方法会高估总风险
    \end{itemize}

    \textbf{B选项错误。}
    这个说法不准确：
    \begin{itemize}
        \item 风险类别间的复合效应在统一方法中更好捕捉
        \item 分块方法可能忽略这些效应
    \end{itemize}

    \textbf{C选项错误。}
    这个说法不准确：
    \begin{itemize}
        \item 这是分块方法的特点
        \item 统一方法考虑风险相关性
    \end{itemize}

    \textbf{D选项正确。}
    这个说法正确：
    \begin{itemize}
        \item 统一方法能更好捕捉错误方向风险
        \item 考虑风险间的相互影响
        \item 通常产生更高的总风险估计
    \end{itemize}
\end{explanation}

\checklist{风险度量方法的比较}




\begin{question}
    The CRO of a small bank has recommended that the bank develop an economic capital framework. In a meeting with the board of directors, the CRO lists several potential benefits to the bank from adopting the framework. Which of the following correctly describes a potential benefit of implementing an economic capital framework?
\end{question}

\begin{choices}
    \item It increases the pricing power of banks by allowing them to charge more competitive interest rates on wholesale loans and retail deposits
    \item It can be extended down to the level of a single customer in order to assess the customer's profitability on a risk-adjusted basis
    \item It can incentivize senior managers by increasing their compensation as the level of economic capital allocated to their business unit increases
    \item It allocates capital based on the portfolio's standalone risk rather than its marginal contribution to the bank's total risk
\end{choices}

\begin{correctanswer}
    B
\end{correctanswer}

\begin{explanation}
    \textbf{A选项错误。}
    这个说法不准确：
    \begin{itemize}
        \item 经济资本框架不直接提高定价能力
        \item 利率定价受市场因素影响
    \end{itemize}

    \textbf{B选项正确。}
    这个说法正确：
    \begin{itemize}
        \item 经济资本可应用于单个客户层面
        \item 用于评估风险调整后的客户盈利性
        \item 有助于客户关系管理
    \end{itemize}

    \textbf{C选项错误。}
    这个说法不准确：
    \begin{itemize}
        \item 经济资本增加不应导致薪酬增加
        \item 应基于风险调整后的收益
    \end{itemize}

    \textbf{D选项错误。}
    这个说法不准确：
    \begin{itemize}
        \item 经济资本应考虑边际贡献
        \item 不应仅基于独立风险
    \end{itemize}
\end{explanation}

\checklist{经济资本框架的实施效益}



\newpage



\section{考点4:大型银行控股公司的资本规划}
\setcounter{question}{0}
\begin{question}
    The Basel Committee recommends that banks use a set of early warning indicators in order to identify emerging risks and potential vulnerabilities in its liquidity position. Which of the following are not early warning indicators of a potential liquidity problem?
    i. Rapid asset growth
    ii. Negative publicity
    iii. Credit rating downgrade
    iv. Increased asset diversification
\end{question}

\begin{choices}
    \item ii and iii
    \item iv only
    \item i and iv
    \item i, ii and iv
\end{choices}

\begin{correctanswer}
    B
\end{correctanswer}

\begin{explanation}
    \textbf{A选项错误。}
    这个说法不准确：
    \begin{itemize}
        \item 负面宣传是预警指标
        \item 信用评级下调是预警指标
    \end{itemize}

    \textbf{B选项正确。}
    这个说法正确：
    \begin{itemize}
        \item 资产多样化不是流动性问题预警指标
        \item 其他三项都是预警指标
    \end{itemize}

    \textbf{C选项错误。}
    这个说法不准确：
    \begin{itemize}
        \item 快速资产增长是预警指标
        \item 资产多样化不是预警指标
    \end{itemize}

    \textbf{D选项错误。}
    这个说法不准确：
    \begin{itemize}
        \item 快速资产增长是预警指标
        \item 负面宣传是预警指标
        \item 资产多样化不是预警指标
    \end{itemize}
\end{explanation}

\checklist{流动性风险的早期预警指标}



\newpage

\end{document}